\section{Tableau Insights}

\subsection{Dashboard layout strategy}
The dashboard is structured for fast escalation from overview to incident details:
\begin{itemize}
	\item KPI summary (total, consensus, anomaly rate),
	\item region-method heatmap and trend slices,
	\item method comparison across the four detectors,
	\item row-level drill-down with method votes and consensus.
\end{itemize}

The layout supports two user personas. Executives use KPI and trend panels to assess system-level risk posture, while analysts and operators use drill-down views to inspect specific timestamps, regions, and detector agreement patterns.

\subsection{Business value for NSPI}
This layout supports faster reliability decisions by helping teams:
\begin{itemize}
	\item prioritize high-volatility regions,
	\item track peak/weather stress concentration,
	\item separate weak alerts from high-confidence consensus events.
\end{itemize}

In practice, this improves prioritization discipline. Incidents with multi-method agreement and high local concentration can be elevated quickly, while isolated single-detector events can be monitored before costly intervention.

\subsection{Explainable AI tooltips}
Tooltips expose triggering methods, vote counts, and anomaly context so operators can quickly understand \emph{why} a point was flagged.

\subsection{Published data products}
The Tableau pipeline uses:
\begin{itemize}
	\item \texttt{anomaly\_results.csv}
	\item \texttt{autoencoder\_results.csv}
	\item \texttt{region\_method\_summary.csv}
	\item \texttt{global\_method\_summary.csv}
\end{itemize}

These outputs support both incident triage and executive monitoring.

The split between row-level and aggregate products is intentional: row-level outputs support forensic analysis and reproducibility, while aggregate products provide concise trend communication for recurring operational reviews.

Additional drill-down visuals are in the appendix: \autoref{fig:app_customer_vs_consumption}, \autoref{fig:app_hour_vs_consumption}, \autoref{fig:app_temp_vs_consumption}, and \autoref{fig:app_pair_plot}.


These outputs support both incident triage and executive monitoring.
